\subsection{General sparsification theorems}
\label{subsec:general-sparsification}

We choose thresholds proportional to factor sizes. Bounding the
number of factors containing each vertex bounds both the total
completion loss and the degree after replacement, provided the
retained edges also have bounded degree.

Use the edge decomposition in~\eqref{eq:factor-edge-union} and the
replacement graph $J$ from Section~\ref{sec:bounded-degree-hardness}.
The input supplies $E_0$, the factor index sets, and each factor's
support or sides. Put $N=|V(H)|$.
Fix constants $d_0\ge0$,
$d_1\ge1$, and assume:
\begin{enumerate}
\item\label{hyp:retained-degree} $\Delta((V(H),E_0))\le d_0$.
\item\label{hyp:factor-incidence}
$|\{f\in\mathcal C\cup\mathcal B:v\in U_f\}|\le d_1$
for every $v\in V(H)$.
\end{enumerate}
Counting incidences between vertices and factor supports gives
\begin{equation}\label{eq:general-factor-incidence}
 \sum_{f\in\mathcal C\cup\mathcal B}|U_f|
 \overset{\text{Hyp.~\ref{hyp:factor-incidence}}}{\leq}d_1N.
\end{equation}

For a fixed real $\zeta>0$, the integer
$t_f=\lceil\zeta|U_f|\rceil\ge1$ is the number of vertices in
each set that the replacement of factor $f$ must join. Write $R_f$
for that replacement graph, and let
$\mathcal R=(R_f)_{f\in\mathcal C\cup\mathcal B}$ be the collection
of all replacements. Each $R_f$ is built on its prescribed support;
the output keeps $E_0$ and takes their union:
\[
 H\xrightarrow{\text{Lemma~\ref{lem:factor-replacement}}}(E_0,\mathcal R)
 \xrightarrow{\text{Lemma~\ref{lem:factor-replacement}}}J,\qquad
 V(J)=V(H),\qquad
 E(J)=E_0\cup\bigcup_{f\in\mathcal C\cup\mathcal B}E(R_f).
\]
The following lemma bounds the degree of this union and the sum of
the threshold losses.

\begin{lemma}[Common sparse factor replacements]
\label{lem:factor-replacement}
Let $H$ have the edge description~\eqref{eq:factor-edge-union},
and let fixed constants $d_0\ge0$, $d_1\ge1$ satisfy
conditions~\ref{hyp:retained-degree}--\ref{hyp:factor-incidence}.
Fix $\zeta>0$ and set
$t_f=\lceil\zeta|U_f|\rceil$ for every factor. There is a
deterministic polynomial-time construction of a spanning subgraph
$J\subseteq H$ with $E_0\subseteq E(J)$ that replaces each factor by a graph
satisfying its no-hole condition at threshold $t_f$, and has maximum degree at most
$d_0+d_1c_\zeta$, where $c_\zeta$ depends only on $\zeta$.
Moreover,
$\sum_{f\in\mathcal C\cup\mathcal B}(t_f-1)<\zeta d_1N$
whenever at least one factor is present.
\end{lemma}

\begin{proof}
Set $\kappa=\max\{1,1/\zeta-1\}$ and
$c_\zeta=d_{\zeta,\kappa}$ from Lemma~\ref{lem:factor-expanders}.
Use that lemma for each clique factor. For a biclique with side
sizes $a,b$, put $u=\lceil\zeta(a+b)\rceil$. If $\min\{a,b\}<u$,
use the empty graph on its support: no admissible pair of $u$-sets
exists. Otherwise $\min\{a,b\}\ge u\ge\zeta(a+b)$, so
\[
 \frac{\max\{a,b\}}{\min\{a,b\}}
 =\frac{a+b}{\min\{a,b\}}-1
 \le\frac1\zeta-1\le\kappa.
\]
Lemma~\ref{lem:factor-expanders} therefore applies to this factor.
Each replacement has degree at most $c_\zeta$ and satisfies the
no-hole condition at threshold $t_f$. If $\zeta>1/2$ the nonempty
branch is impossible, and if $\zeta=1/2$
it requires $a=b$.

Take the union of the replacements with $E_0$. The replacements of
vacuous biclique factors need not be connected. The completion
lemmas require only the no-hole conditions; when connected factor
replacements are needed separately, use Lemma~\ref{lem:factor-expanders}
with its side-ratio assumption.
For every $v\in V(H)$, the degree in the union satisfies
\[
 \deg_J(v)\overset{\text{Lem.~\ref{lem:factor-expanders}}}{\leq}
 \deg_{(V(H),E_0)}(v)
 +|\{f\in\mathcal C\cup\mathcal B:v\in U_f\}|c_\zeta
 \overset{\text{Hyps.~\ref{hyp:retained-degree}, \ref{hyp:factor-incidence}}}{\leq}
 d_0+d_1c_\zeta.
\]
Finally, the ceiling defining $t_f$ gives
\begin{equation}\label{eq:factor-threshold-slack}
 t_f-1<\zeta|U_f|\qquad(f\in\mathcal C\cup\mathcal B).
\end{equation}
If at least one factor is present, then
\[
 \sum_{f\in\mathcal C\cup\mathcal B}(t_f-1)
 \overset{\eqref{eq:factor-threshold-slack}}{<}
 \zeta\sum_{f\in\mathcal C\cup\mathcal B}|U_f|
 \overset{\eqref{eq:general-factor-incidence}}{\leq}\zeta d_1N.
\]
Each factor has nonempty support, so there are at most
$d_1N$ factors. Let $h\ge1$ be the fixed walk length chosen in
Appendix~\ref{app:factor-expanders} for this $\zeta$. The construction
enumerates $O(8^h\sum_{f\in\mathcal C\cup\mathcal B}|U_f|)
=O(8^h d_1N)$ auxiliary steps.
Sorting the resulting endpoint pairs and merging repetitions uses
$O(8^h d_1N\log(2+8^h d_1N))$ comparisons, so the construction
is deterministic and polynomial time for the fixed parameters.
\end{proof}

\begin{theorem}[Treedepth sparsification]
\label{thm:td-sparsification}
Let $H$ have the edge description~\eqref{eq:factor-edge-union},
and let fixed constants $d_0\ge0$, $d_1\ge1$ satisfy
conditions~\ref{hyp:retained-degree}--\ref{hyp:factor-incidence}.
For every fixed $\varepsilon>0$
one can construct deterministically in polynomial time a spanning
subgraph $J\subseteq H$ with $E_0\subseteq E(J)$ whose maximum degree depends only on
$d_0,d_1,\varepsilon$ and such that
$\td(H)-\varepsilon N\le\td(J)\le\td(H)$.
Every elimination forest of $J$ can be transformed into one of
$H$ whose height increases by at most $\varepsilon N$.
\end{theorem}

\begin{proof}
Apply Lemma~\ref{lem:factor-replacement} with
$\zeta=\varepsilon/(4d_1)$. Its degree bound depends only on
the stated constants. The exception bound in
Lemma~\ref{lem:factor-sparsification} satisfies
\begin{equation}\label{eq:common-forest-exception-budget}
 \operatorname{exc}_{\mathrm{td}}
 \overset{\eqref{eq:completion-deletion-bounds}}{=}
 2\sum_{f\in\mathcal C}(t_f-1)+4\sum_{f\in\mathcal B}(t_f-1) \le4\sum_{f\in\mathcal C\cup\mathcal B}(t_f-1)
 \overset{\text{Lem.~\ref{lem:factor-replacement}}}{\leq}
 4\zeta d_1N=\varepsilon N.
\end{equation}
Lemma~\ref{lem:factor-sparsification} now gives both conclusions.
\end{proof}

\begin{theorem}[Treewidth sparsification]
\label{thm:tw-sparsification}
Let $H$ have the edge description~\eqref{eq:factor-edge-union},
and let fixed constants $d_0\ge0$, $d_1\ge1$ satisfy
conditions~\ref{hyp:retained-degree}--\ref{hyp:factor-incidence}.
Fix an integer $s\ge1$. Suppose the supplied factor description also
includes, for every $f\in\mathcal B$, integers
$1\le\ell_f,m_f\le s$ and partitions
$A_f=A_{f,1}\mathbin{\dot\cup}\cdots\mathbin{\dot\cup}A_{f,\ell_f}$
and $B_f=B_{f,1}\mathbin{\dot\cup}\cdots\mathbin{\dot\cup}B_{f,m_f}$
into nonempty sets, together with indices
$g_{f,i},h_{f,j}\in\mathcal C$ such that
$A_{f,i}\subseteq U_{g_{f,i}}$ and $B_{f,j}\subseteq U_{h_{f,j}}$
for every $i\in[\ell_f]$ and $j\in[m_f]$.
For every fixed $\varepsilon>0$, one can construct deterministically
in polynomial time a spanning subgraph $J\subseteq H$ with
$E_0\subseteq E(J)$ whose maximum
degree depends only on $d_0,d_1,s,\varepsilon$ and such that
$\operatorname{tw}(H)-\varepsilon N
 \le\operatorname{tw}(J)\le\operatorname{tw}(H)$.
Every tree decomposition of $J$ can be transformed into one of
$H$ whose width increases by at most $\varepsilon N$.
\end{theorem}

\begin{proof}
Apply Lemma~\ref{lem:factor-replacement} to the original factors
with $\zeta=\varepsilon/(d_1\max\{3,2s^2\})$. It constructs $J$
with maximum degree depending only on the stated constants and
thresholds $t_f=\lceil\zeta|U_f|\rceil$.

For $f\in\mathcal B$, the supplied partitions and their indices give
\begin{equation}\label{eq:tw-factor-piece-count}
 \ell_fm_f\leq s^2,\qquad
 A_f\times B_f=\mathop{\dot\bigcup}_{i\in[\ell_f],\,j\in[m_f]}
 (A_{f,i}\times B_{f,j}).
\end{equation}
Describe the same graph $H$ by keeping $E_0$ and its clique
factors and replacing the biclique indexed by $f$ with the
$\ell_fm_f$ bicliques indexed by $(f,i,j)$, on sides
$A_{f,i},B_{f,j}$. Give each of these bicliques threshold $t_f$.
Each has its sides contained in the clique supports $U_{g_{f,i}}$
and $U_{h_{f,j}}$.

The same graph $J$ satisfies this refined factor description. The
clique conditions hold by Lemma~\ref{lem:factor-replacement}.
Any $t_f$-subsets of $A_{f,i}$ and $B_{f,j}$ are also subsets of
$A_f$ and $B_f$, so the original replacement joins them. If either
new side has fewer than $t_f$ vertices, this condition is vacuous.

Lemma~\ref{lem:tw-factor-sparsification} charges at most $3(t_f-1)$
per clique and $2(t_f-1)$ per refined biclique. Thus its completion
increases width by at most
\[
\begin{aligned}
 &3\sum_{f\in\mathcal C}(t_f-1)
 +2\sum_{f\in\mathcal B}\ell_fm_f(t_f-1)
 \quad\overset{\eqref{eq:tw-factor-piece-count}}{\leq}
 3\sum_{f\in\mathcal C}(t_f-1)
       +2s^2\sum_{f\in\mathcal B}(t_f-1)\\
 &\quad\le\max\{3,2s^2\}\sum_{f\in\mathcal C\cup\mathcal B}(t_f-1)
 \quad\overset{\text{Lem.~\ref{lem:factor-replacement}}}{\leq}
 \max\{3,2s^2\}\zeta d_1N=\varepsilon N.
\end{aligned}
\]
The comparison lemma also gives the stated parameter inequalities.
\end{proof}

\Needspace{9\baselineskip}
\begin{theorem}[Total-depth sparsification]
\label{thm:total-depth-sparsification}
Let $H$ have the edge description~\eqref{eq:factor-edge-union},
and let fixed constants $d_0\ge0$, $d_1\ge1$ satisfy
conditions~\ref{hyp:retained-degree}--\ref{hyp:factor-incidence}.
For every fixed $\varepsilon>0$
one can construct deterministically in polynomial time a spanning
subgraph $J\subseteq H$ with $E_0\subseteq E(J)$ whose maximum degree depends only on
$d_0,d_1,\varepsilon$ and such that
$\totd(H)-\varepsilon N^2\le\totd(J)\le\totd(H)$.
Every elimination forest $F$ of $J$ can be transformed into an
elimination forest $\widehat F$ of $H$ satisfying
$\totd(\widehat F)\le\totd(F)+\varepsilon N^2$.
\end{theorem}

\begin{proof}
Use the same factor replacements as in
Theorem~\ref{thm:td-sparsification}, with
$\zeta=\varepsilon/(4d_1)$. Equation~\eqref{eq:common-forest-exception-budget}
gives $\operatorname{exc}_{\mathrm{td}}\le\varepsilon N$, and
Lemma~\ref{lem:total-depth-comparison} gives the completion and both
parameter bounds, with total-depth increase at most
$N\operatorname{exc}_{\mathrm{td}}
\overset{\eqref{eq:common-forest-exception-budget}}{\leq}\varepsilon N^2$.
\end{proof}

If a construction has a fixed number $d$ of designated vertices
that are not needed for sparsification, deleting them lowers treedepth and treewidth by
at most $d$. Restoring them above all existing roots increases total
depth by at most $dN$, where $N$ is the original graph order.
These costs are $O(1)$ for treedepth and treewidth, whose promise
bounds are linear in $N$, and $O(N)$ for total depth, whose bounds
are quadratic. They are therefore $o(Y)$ for the corresponding YES
upper bound $Y$, so
Lemma~\ref{lem:gap-preservation} absorbs them on sufficiently large
instances. We use this to omit the three vertices
$z_A,z_B,z_C$ from the NP treedepth construction.
